\chapter{Die Modellbibliothek}
\label{kap:editor-bibliothek}

Halbleiter brauchen Modellparameter, und fertige Bausteine wie
Operationsverstärker kommen als Teilschaltungen. Die Modellbibliothek
sammelt solche Modelle aus mehreren Quellen und kopiert ein gewähltes Modell
in die Datei. Die Datei bleibt dadurch vollständig: Sie rechnet auch auf
einem Rechner, der die Bibliothek nicht kennt.

\section{Quellen}

\begin{description}[style=nextline,leftmargin=2em]
  \item[Standardbibliothek]
    Ins Programm eingebaut (\datei{editor/library/standard.lib}). Sie
    enthält nur allgemeine Modelle mit Lehrbuchwerten, keine Bauteile
    bestimmter Hersteller:

    \begin{tabularx}{\linewidth}{@{}lL@{}}
      \toprule
      Name & Beschreibung \\
      \midrule
      \code{NPN_GEN}, \code{PNP_GEN} & Allgemeiner NPN- und PNP-Transistor \\
      \code{D_GEN} & Allgemeine Siliziumdiode, etwa 0,6\,V bei 1\,mA \\
      \code{Z5V1_GEN} & Z-Diode 5,1\,V, vereinfacht \\
      \code{NMOS_L1}, \code{PMOS_L1} & MOSFETs Level 1 mit Gate-Ladung \\
      \code{n18}, \code{p18} & MOSFETs BSIM3v3, typisch für 0,18\,µm \\
      \bottomrule
    \end{tabularx}
  \item[Beispiele]
    Die Modelle aus \code{"models"} der mitgelieferten Beispiele.
  \item[Eigene Bibliotheken]
    Dateien und ganze Ordner, die man selbst einbindet:
    \begin{itemize}
      \item SPICE-Dateien (\datei{.lib}, \datei{.mod}, \datei{.txt},
        \datei{.cir}, \datei{.sp}) mit \code{.model}-Karten und
        \code{.subckt}-Teilschaltungen, wie Hersteller sie veröffentlichen,
      \item esim-Dateien (\datei{.json}) mit ihrem Abschnitt
        \code{"models"}.
    \end{itemize}
\end{description}

\section{Eigene Bibliotheken einbinden}

Im Dialog \menue{Modellbibliothek} (\tasten{Strg}{L}) führt
\menue{Bibliotheken \dots} zur Liste der eigenen Bibliotheken.
\menue{Datei \dots} und \menue{Ordner \dots} fügen hinzu,
\menue{Entfernen} nimmt einen Eintrag wieder heraus. Ein Ordner wird mit
allen passenden Dateien gelesen. Die Liste bleibt über den Neustart
erhalten.

Beim Lesen einer SPICE-Datei gilt:
\begin{itemize}
  \item Fortsetzungszeilen (\code{+}) und Kommentare (\code{*}, \code{;})
    werden verstanden.
  \item Jede \code{.model}-Karte auf oberster Ebene ist ein Modell. Der
    Kommentar direkt über der Karte ist ihre Beschreibung; dessen erste
    Zeile erscheint in der Liste.
  \item Jeder \code{.subckt}-Block ist eine Teilschaltung. Er nimmt die
    \code{.model}- und \code{.subckt}-Definitionen der Datei mit, die er
    benutzt. Reine Hilfs-Teilschaltungen, die nur von anderen benutzt
    werden, erscheinen nicht in der Liste. Enthält eine Datei nur eine
    Teilschaltung, ist der Kommentarkopf der Datei ihre Beschreibung.
  \item Modelltypen, die esim nicht kennt, werden übergangen; was sich
    nicht lesen lässt, nennt der Dialog unter der Liste.
\end{itemize}

\begin{hinweis}[Modelle von Herstellern]
  Die Modelldateien der Halbleiterhersteller stehen meist auf deren
  Produktseiten (»SPICE model«, »PSpice model«). esim versteht die
  üblichen PSpice-Formen (Kapitel~\ref{kap:format-teilschaltungen}) und
  ist mit Makromodellen von Operationsverstärkern und mit
  Leistungs-MOSFETs (Level~3) erprobt. Weil die Lizenz solcher Dateien
  die Weitergabe oft nicht erlaubt, gehören sie in die eigene Bibliothek,
  nicht in ein öffentliches Projekt.
\end{hinweis}

\section{Der Dialog}

\menue{Bearbeiten}\pfad\menue{Modellbibliothek \dots} (\tasten{Strg}{L})
oder der Eintrag \menue{Modellbibliothek \dots} der Bauteilliste öffnet
den Dialog:

\begin{itemize}
  \item \textbf{Suchfeld}: sucht in Name, Beschreibung und Quelle.
  \item \textbf{Art}: Alle Arten, Dioden, Bipolartransistoren, MOSFETs,
    JFETs oder Teilschaltungen (SPICE).
  \item \textbf{Liste} mit Name, Art, Beschreibung und Quelle.
  \item \textbf{Karte}: der vollständige Text des gewählten Modells mit
    seiner Quelle.
\end{itemize}

Die Schaltflächen unten:
\begin{description}[leftmargin=2em]
  \item[Im Plan platzieren] kopiert das Modell in die Datei und hängt ein
    passendes Bauteil (Diode, Transistor, Instanz der Teilschaltung) mit
    diesem Modell an den Mauszeiger.
  \item[In die Datei übernehmen] kopiert nur das Modell nach
    \code{"models"} bzw. die Teilschaltung nach \code{"subcircuits"}.
  \item[Verwenden] (wenn der Dialog aus dem Feld \code{model} der
    Eigenschaften geöffnet wurde) setzt das Modell am gewählten Bauteil.
\end{description}

\section{Schneller: Suche in der Bauteilliste}

Für den Alltag ist der Dialog oft nicht nötig. Die Suche der Bauteilliste
findet auch die Modelle der Bibliothek: \taste{A}, »NPN\_GEN« tippen,
Eingabetaste – der Transistor hängt mit seinem Modell am Mauszeiger. Ebenso
schlägt das Feld \code{model} in den Eigenschaften beim Tippen passende
Modelle vor.

\section{Was in der Datei landet}

Ein Modell wird als SPICE-Zeile in \code{"models"} kopiert, lange Karten als
Liste ihrer Zeilen:

\begin{json}
"models": {
  "NPN_GEN": ".model NPN_GEN NPN(IS=1e-14 BF=200 BR=2 CJE=10p CJC=5p TF=0.5n)",
  "n18": [".model n18 NMOS (LEVEL=49 VERSION=3.1 TNOM=27 TOX=4.1E-9",
          "+ XJ=1E-7 NCH=2.3549E17 VTH0=0.3694303 ...)"]
}
\end{json}

Eine Teilschaltung landet als SPICE-Text unter \code{"subcircuits"},
zusammen mit allem, was sie braucht. Ist ein Modell gleichen Namens schon in
der Datei, wird es nicht noch einmal kopiert.
